\section*{摘要}

随着2022年FIA引入"地面效应"空气动力学规则，
F1赛车的海豚跳（Porpoising）现象成为工程设计中的关键挑战。
传统CFD仿真在规则限制下难以高效覆盖大规模参数空间，
亟需低成本、高效率的替代方法。

本文基于Kaggle公开的F1气动稳定性数据集（约15万样本），
构建了一套数据驱动的代理优化框架。
首先通过6种机器学习模型的系统对比，
选定XGBoost与DeepEnsemble作为代理模型，
分别实现R$^2 = 0.9999$和R$^2 = 0.9995$的预测精度。
在此基础上，
提出风险敏感粒子群优化算法（Risk-Sensitive PSO），
以$\mu(x)-\lambda\sigma(x)$为适应度函数，
有效抑制代理模型在分布外（OOD）区域的幻觉行为。
进一步通过主动学习闭环实现模型精炼，
将MSE从0.3272降至0.1444（降低55.9\%），
收敛速度提升2.1倍。

在多场景优化模块中，
针对Monza高速、Monaco高下压力、均衡设置和湿地四种典型F1工况，
采用多目标加权适应度函数进行Pareto前沿搜索，
并通过SHAP、排列重要性、约束灵敏度分析和Landscape Diagnosis
对最优解进行了系统性的可解释性验证。
Porpoising风险分析模块通过一阶梯度热力图和二阶Hessian曲率分析，
刻画了不同工况下的海豚跳风险空间分布。

实验结果表明：
代理模型R$^2$达到0.9999，远超理想目标（$\geq 0.92$）；
PSO平均收敛代数为27.9代（最优目标$\leq 50$）；
单次寻优耗时约83ms；
消融实验验证了风险敏感适应度的OOD保护功能，
以及多目标优化相较单目标优化的区分度提升（从0.1\%到20\%以上）。

研究工作形成了一套完整的"数据预处理—代理建模—风险感知优化—可解释性验证"
技术路线，
为F1赛车气动参数的高效求解提供了可量化的方法论支撑。

\vspace{0.5cm}

\noindent
\textbf{关键词：}
代理模型；
粒子群优化；
F1赛车；
气动稳定性；
海豚跳；
多目标优化；
可解释性分析
